
%==========================================================================
\section{Analyse expérimentale}\label{sec:exp}
%==========================================================================

Chaque chiffre ci-dessous est produit par \texttt{src/experiments.py}, stocké en CSV
brut dans \texttt{results/}, puis transformé en figure par \texttt{src/plots.py} et en
tableau LaTeX par \texttt{report/tables.py}. Aucun chiffre de ce rapport n'est saisi à
la main. Le matériel et les deux pièges de configuration sont décrits en
Section~\ref{sec:setup} ; sauf mention contraire $K=32$, $\lambda=0,01$, $c_0=512$,
$\alpha=0,4$, graine 42, et 3 répétitions de 4 itérations dont la première est
écartée comme échauffement. La ligne de commande exacte de chaque exécution figure
dans \texttt{scripts/run\_all.sh}.

\paragraph{Avertissement sur la comparabilité des mesures.} Toute la campagne est une
séquence de jobs à 8 cœurs sur un ordinateur portable, et un portable
sous charge soutenue ne maintient pas une horloge constante. Nous avons mesuré la
conséquence plutôt que de la balayer d'un revers de main : la même configuration
ré-exécutée une heure plus tard, après un voisin plus lourd, est ressortie jusqu'à
deux fois plus lente. Nous attribuons cela à une combinaison de limitation thermique
(\emph{throttling}) et de workers JVM/Python froids dans chaque \texttt{SparkSession}
nouvellement créée, et nous n'avons pas cherché à séparer les deux effets. La règle
pratique que nous avons suivie, et que le lecteur devrait appliquer en lisant les
tableaux : \textbf{chaque conclusion ci-dessous est tirée d'un rapport calculé
\emph{à l'intérieur} d'un seul appel d'expérience} - une accélération, un facteur de
croissance, un rapport entre méthodes - et jamais de secondes absolues comparées
entre deux appels. Là où une expérience aurait sinon été exposée à cette dérive, elle
a été ré-exécutée avec les répétitions entrelacées entre configurations et un temps de
repos entre elles (\texttt{--interleave}, \texttt{--cooldown}).

%--------------------------
\subsection{Exactitude de la parallélisation}\label{sec:correctness}
%--------------------------

Rien d'autre dans cette section n'a de sens si le code distribué ne calcule pas eALS.
Deux vérifications sont effectuées, dans cet ordre. D'abord, les règles de mise à jour
elles-mêmes sont comparées à une évaluation littérale en $O(MNK)$ de
l'Éq.~\eqref{eq:eq5} sur une matrice jouet, qui concorde à $3\cdot10^{-16}$ près
(Section~\ref{sec:algo}). Ensuite - le tableau ci-dessous - l'implémentation
distribuée est comparée à la référence séquentielle. La Section~\ref{sec:algo} a
argumenté que le découpage en blocs ne peut pas changer le résultat ; ici nous le
vérifions. L'implémentation RDD est exécutée sur ml-100k avec la même graine, $K=8$ et
5 itérations, pour 1, 2, 4 et 8 blocs, et comparée entrée par entrée à la référence
NumPy séquentielle.

\begin{table}[htbp]
\centering\small
\caption{Contrôle d'exactitude (\texttt{check\_correctness.py}, ml-100k, $K=8$, 5
itérations). « blocs » est le nombre de partitions Spark. L'objectif est
l'Éq.~\eqref{eq:obj7} évaluée via l'Éq.~\eqref{eq:objfast}.}\label{tab:correctness}
\begin{tabular}{@{}lrrrrr@{}}
\toprule
Implémentation & blocs & $\max|\Delta P|$ & $\max|\Delta Q|$ & Objectif $J$ & err.\ rel. \\
\midrule
\inputrows{tables/tab_correctness.tex}
\bottomrule
\end{tabular}
\end{table}

Deux observations. Avec un seul bloc, l'exécution Spark est \textbf{identique au bit
près} à la référence NumPy - la même arithmétique dans le même ordre. Avec 2, 4 ou 8
blocs, l'écart est de $\CorrWorst$, soit quelques unités du dernier chiffre d'un
double : la seule chose qui change est l'ordre dans lequel les sommes partielles de
$S^p$ et $S^q$ sont accumulées. L'objectif concorde à une erreur relative inférieure à
$10^{-15}$. C'est la forme empirique de l'affirmation de l'article selon laquelle eALS
est trivialement parallélisable \emph{sans perte d'approximation}, et cela signifie
aussi que chaque mesure de temps ci-dessous mesure le même calcul.

%--------------------------
\subsection{Scalabilité}\label{sec:scal}
%--------------------------

\subsubsection{Scalabilité forte}\label{sec:strong}

Données fixes, ressources croissantes : \texttt{local[1]}, \texttt{local[2]},
\texttt{local[4]}, \texttt{local[8]}, avec un nombre de partitions maintenu
proportionnel au nombre de cœurs pour que le travail par tâche reste constant.

\begin{figure}[htbp]\centering
\figplot{strong_scaling}{0.98\textwidth}
\caption{Scalabilité forte sur ml-1m ($10^6$ interactions) et ml-10m ($10^7$). Les
barres d'erreur sont le min--max sur 9 échantillons par itération. Pointillés :
l'accélération linéaire idéale ; tirets : la courbe d'Amdahl ajustée à chaque jeu de
données.}\label{fig:strong}
\end{figure}

\begin{table}[htbp]
\centering\small
\caption{Scalabilité forte, $K=32$, partitions $=2\times$ cœurs.}\label{tab:strong}
\begin{tabular}{@{}lrrrr@{}}
\toprule
Cœurs $p$ & médiane s/itér. & min--max & accélération $S(p)$ & efficacité $E(p)$ \\
\midrule
\inputrows{tables/tab_strong.tex}
\bottomrule
\end{tabular}
\end{table}

\paragraph{Méthodologie.} Le nombre de blocs croît avec le nombre de cœurs ($B=2p$),
de sorte que le travail par tâche reste constant - ce que l'on ferait en production.
La contrepartie, qui compte pour la lecture des chiffres, est que la surcharge totale
par tâche croît elle aussi avec $p$ : à $p=8$ l'exécution paie seize surcharges de
bloc contre deux à $p=1$, ce qui \emph{sous-estime} l'accélération. Les répétitions
sont \emph{entrelacées} - la boucle externe est la répétition, la boucle interne le
nombre de cœurs - afin que la lente dérive thermique d'un portable sous une campagne
prolongée ne retombe pas entièrement sur la configuration exécutée en dernier. Nous
avons constaté que cela comptait : dans une première exécution non entrelacée, le point
à 8 cœurs a produit un unique échantillon de 19,7\,s contre une médiane de 8,8\,s.

\paragraph{Interprétation de l'ajustement d'Amdahl.} En suivant Amdahl~\cite{amdahl1967},
nous ajustons $T(p)=a+b/p$ par moindres carrés et rapportons $s=a/(a+b)$, la fraction
d'une itération à un cœur qui ne bénéficie pas de cœurs supplémentaires. Ajuster le
modèle directement est plus robuste que de lire $s$ sur le point à 8 cœurs, qui est
exactement le point le plus exposé à la contention mémoire.

\paragraph{Résultats.} L'exécution de référence - ml-10m, répétitions
entrelacées - atteint $S(8)=\SpeedupBig\times$ sur un possible $8\times$, une
efficacité de $26\,\%$, avec une fraction série d'Amdahl ajustée de $\SerialBig$. Sur
ml-1m, où les données sont dix fois plus petites et où le coût fixe d'une itération en
représente une part bien plus grande, $S(8)=\SpeedupSmall\times$ et $s=\SerialSmall$.
Les deux courbes s'aplatissent entre 4 et 8 cœurs : le quatrième cœur rapporte encore
quelque chose, le huitième pratiquement pas.

\paragraph{Origine de la fraction non parallélisable.} Quatre choses ne se réduisent pas quand $p$
croît :
\begin{itemize}\itemsep2pt
\item \textbf{La désérialisation des broadcasts}, payée une fois par worker Python et
par broadcast, et donc \emph{croissante} avec $p$. C'est le terme dominant, et il est
invisible dans toute décomposition mesurée côté driver : \texttt{sc.broadcast} n'écrit
le tableau qu'une fois, tandis que la moitié coûteuse - la désérialisation - se
produit dans chaque worker Python et est imputée à la tâche, pas au driver.
\item \textbf{La colle côté driver} : créer les broadcasts, disperser $2\times$ $B$
blocs dans $P$ et $Q$, réduire les partiels de Gram.
\item \textbf{La soumission des jobs et l'ordonnancement des tâches}, $2B$ tâches par
itération.
\item \textbf{La bande passante mémoire}, qui n'est pas du tout du travail série mais
se comporte comme tel : le noyau est un flux de gathers et de sommes par segment sur
des tableaux bien plus grands que le L2, et les huit cœurs d'un portable série M
partagent un seul contrôleur mémoire. C'est pourquoi le plafond est une propriété de
la machine plutôt que du code, et pourquoi un cluster de huit nœuds à un cœur
paraîtrait meilleur qu'un seul nœud à huit cœurs.
\end{itemize}

\subsubsection{Scalabilité en nombre de facteurs latents}\label{sec:factors}

C'est l'expérience qui décide si l'affirmation centrale de l'article survit au contact
d'une implémentation réelle. Par itération, eALS coûte
$O\big((M{+}N)K^2+|\mathcal{R}|K\big)$ et ALS $O\big((M{+}N)K^3+|\mathcal{R}|K^2\big)$.
Lire cela comme « $K^2$ contre $K^3$ » est une erreur, et l'erreur se manifeste
immédiatement dans les mesures. Les deux bornes sont la \emph{somme} de deux termes,
et lequel domine dépend de la forme des données :
\begin{equation}\label{eq:kstar}
(M{+}N)K^2 \;\ge\; 2|\mathcal{R}|K
\quad\Longleftrightarrow\quad
K \;\ge\; K^{\star}=\frac{2|\mathcal{R}|}{M+N}.
\end{equation}
Sous $K^{\star}$, eALS est essentiellement \emph{linéaire} en $K$ ; ce n'est
qu'au-dessus que le terme quadratique prend le dessus. $K^{\star}=\CrossKMl$ sur ml-1m
et $K^{\star}=\CrossKAmz$ sur Amazon Movies, donc les deux jeux de données se situent
de part et d'autre de la frontière intéressante et nous exécutons le balayage sur les
deux. (Pour ALS, le même raisonnement donne un régime cubique au-dessus de
$|\mathcal{R}|/(M{+}N)$, c.-à-d.\ la moitié de $K^{\star}$.)

\begin{figure}[htbp]\centering
\figplot{factor_scaling}{0.92\textwidth}
\caption{Coût d'une itération en fonction de $K$, 8 cœurs, log-log. Les repères gris
sont $K$, $K^2$ et $K^3$ purs, ancrés au plus petit $K$. « plage haute » est la pente
ajustée sur les deux dernières octaves seulement, où le coût constant de Spark ne
compte plus.}\label{fig:factors}
\end{figure}

\begin{table}[htbp]
\centering\small
\caption{Secondes par itération en fonction de $K$, 8 cœurs. eALS est notre
implémentation Spark RDD ; ALS est
\texttt{pyspark.ml.recommendation.ALS(implicitPrefs=True)}, c.-à-d.\ Hu et
al.~\cite{hu2008implicit}, chronométré par son coût \emph{marginal}
$(T_{11}-T_1)/10$ afin d'exclure le coût fixe de construction de ses blocs. ALS n'a pas
été exécuté au-delà de $K=128$ : une seule itération y coûte déjà des
minutes.}\label{tab:factors}
\begin{tabular}{@{}lrrr@{}}
\toprule
$K$ & eALS (s/itér.) & MLlib ALS (s/itér.) & rapport \\
\midrule
\inputrows{tables/tab_factors.tex}
\bottomrule
\end{tabular}
\end{table}

\paragraph{Résultats.} L'affirmation la plus nette est un facteur de
croissance sur une fenêtre fixe plutôt qu'un exposant ajusté. Sur \textbf{Amazon
Movies}, entre $K=32$ et $K=128$, le coût d'une itération eALS est multiplié par
$\GrowthEalsAmz$ et celui d'une itération MLlib ALS par $\GrowthAlsAmz$ - ALS croît
près de cinq fois plus vite sur la même fenêtre, ce qui est l'affirmation qualitative
de l'article. La conséquence en termes absolus est un point de croisement : à
$K=\KAtRatioMinAmz$, MLlib ALS est $\RatioAtMinAmz\times$ \emph{plus rapide} que notre
eALS, et à $K=\KAtRatioAmz$, eALS est $\RatioAtMaxAmz\times$ plus rapide. Sur
\textbf{ml-1m}, la même fenêtre donne $\GrowthEalsMl$ contre $\GrowthAlsMl$ : le
rapport entre les deux taux de croissance est essentiellement le même, sur un jeu de
données dont la forme est complètement différente.

\paragraph{La forme théorique s'ajuste, et elle nous dit où se trouve réellement le
croisement.} Ajuster $t = a + bK + cK^2$ pour eALS et $t = a + bK^2 + cK^3$ pour ALS
- les deux bornes de complexité, avec des constantes libres - reproduit les
mesures presque exactement : $R^2 = \FitRtwoEalsAmz$ et $\FitRtwoAlsAmz$ sur Amazon
Movies, $\FitRtwoEalsMl$ et $\FitRtwoAlsMl$ sur ml-1m. Les constantes additives sont
$\FitConstEalsAmz$\,s pour eALS contre $\FitConstAlsAmz$\,s pour MLlib : notre
implémentation Python paie environ le double du coût fixe par itération du code JVM
natif, ce qui explique qu'ALS gagne à petit $K$.

Le chiffre intéressant est le rapport $b/c$, le $K$ auquel le terme d'ordre supérieur
dépasse le terme inférieur. Le décompte de flops prédit
$K^{\star}=2|\mathcal{R}|/(M{+}N)=\CrossKAmz$ sur Amazon Movies ; le croisement
\emph{ajusté} est $\FitCrossEalsAmz$, neuf fois plus élevé. Sur ml-1m la prédiction
est $\CrossKMl$ et l'ajustement donne $\FitCrossEalsMl$.

\paragraph{Pourquoi le décompte de flops est faux d'un ordre de grandeur, et pourquoi
c'est la chose la plus utile que dit cette expérience.} Les deux termes de
$O((M{+}N)K^2+|\mathcal{R}|K)$ ont des constantes \emph{par opération} très
différentes dans toute implémentation réelle, et la borne de complexité les traite
comme interchangeables. Dans notre noyau :
\begin{itemize}\itemsep2pt
\item le terme en $K^2$ est \texttt{Sq[:, f] @ PT}, un produit matrice-vecteur dense
que BLAS exécute à plusieurs GFLOP/s hors cache ;
\item le terme en $|\mathcal{R}|K$ est \texttt{QT[f][idx]} suivi de deux sommes par
segment \texttt{np.bincount} - un gather aléatoire et un scatter aléatoire sur des
tableaux de $\textit{nnz}$ éléments, limités par la latence mémoire, un ordre de
grandeur plus lent par opération.
\end{itemize}
Un décompte de flops place donc $K^{\star}$ beaucoup trop bas, et l'exposant mesuré
reste proche de $\SlopeHiEalsAmz$ sur le haut de la plage où un décompte de flops pur
prédirait presque $2$. Ce n'est pas une erreur de l'analyse de l'article, qui est
asymptotiquement correcte ; c'est l'écart ordinaire entre un décompte d'opérations et
une hiérarchie mémoire, et cela a une conséquence pratique : \textbf{sur nos deux
jeux de données, eALS se comporte, sur toute la plage de $K$ qu'on peut se permettre
d'entraîner, bien plus près d'un algorithme $O(|\mathcal{R}|K)$ que d'un algorithme
$O((M{+}N)K^2)$} - ce qui est meilleur pour eALS que ne le suggère sa propre borne
de complexité.

\paragraph{Réserve sur la méthode de référence.} L'exposant de MLlib ($\SlopeHiAlsAmz$ sur la
plage haute) est lui aussi bien en dessous de son $3$ nominal, pour une raison du même
type plus une autre : MLlib n'utilise pas une inversion de manuel. Elle accumule une
équation normale et la factorise par Cholesky sur un stockage triangulaire compact, en
code natif. L'article fait la même observation à propos de sa propre référence ALS en
Java. La formulation honnête de cette expérience est donc « un algorithme quadratique
en $K$ écrit en Python contre un algorithme quasi-cubique bien optimisé écrit en code
natif » - et le quadratique gagne quand même à partir de $K=\FlipKAmz$, par une
marge qui croît.

\subsection{Convergence et précision}\label{sec:quality}
%--------------------------

\begin{figure}[htbp]\centering
\figplot{convergence}{0.92\textwidth}
\caption{L'objectif~\eqref{eq:obj7}, calculé via~\eqref{eq:objfast}, à chaque
itération ($K=32$, $c_0=512$, $\alpha=0,4$). La décroissance monotone est la preuve
pratique que les règles de mise à jour sont les minimiseurs de coordonnées
exacts.}\label{fig:conv}
\end{figure}

\begin{table}[htbp]
\centering\small
\caption{Convergence de l'objectif, $K=32$, 8 cœurs. « itérations à 1\,\% » est la
première itération dont l'objectif est à moins de 1\,\% de la valeur après 30
itérations.}\label{tab:convergence}
\begin{tabular}{@{}lrrrrr@{}}
\toprule
Jeu de données & $J$ après 1 itér. & $J$ après 30 & décroissance totale & itérations à 1\,\% & secondes \\
\midrule
\inputrows{tables/tab_convergence.tex}
\bottomrule
\end{tabular}
\end{table}

L'objectif décroît de façon monotone à chaque itération sur les deux jeux de données,
ce qui est le contrôle pratique que~\eqref{eq:pu} et~\eqref{eq:qi} sont réellement les
minimiseurs de coordonnées exacts - un pas de descente de coordonnées même
légèrement faux se manifesterait immédiatement ici par une remontée. La convergence
est rapide dans le sens qui compte : $\ConvItersMl$ itérations sur ml-1m et
$\ConvItersAmz$ sur Amazon Movies amènent $J$ à moins de 1\,\% de sa valeur après 30,
et l'essentiel de la décroissance totale se produit dans les trois premières. Un budget
de 20 itérations, utilisé partout ailleurs dans ce rapport, est donc largement au-delà
du point où l'objectif cesse de bouger.

\begin{table}[htbp]
\centering\small
\caption{Protocole hors-ligne : leave-one-out, classement contre le catalogue
\emph{entier}. $K=32$, $\lambda=0,01$, $c_0=512$, $\alpha=0,4$, 8 cœurs. « fit » est le
temps d'entraînement (horloge murale) du modèle rapporté.}\label{tab:quality}
\begin{tabular}{@{}lrrrrr@{}}
\toprule
Méthode & HR@10 & HR@100 & NDCG@10 & NDCG@100 & fit (s) \\
\midrule
\inputrows{tables/tab_quality.tex}
\bottomrule
\end{tabular}
\end{table}

\begin{figure}[htbp]\centering
\figplot{quality_vs_time}{0.92\textwidth}
\caption{Précision en fonction du temps d'entraînement \emph{réel} (horloge murale),
pas du numéro d'itération. Chaque marqueur est une exécution d'entraînement complète
de la durée donnée, évaluée sur l'ensemble retenu.}\label{fig:qualtime}
\end{figure}

Trois constats, sur Amazon Movies, le jeu de données de l'article.

\textbf{(i) eALS bat nettement l'ALS de MLlib en précision.} HR@100 passe de $\HRa$ à
$\HRe$, soit $+\EalsOverAls$\,\%, avec NDCG@100 s'améliorant dans la même proportion.
Les deux sont entraînés sur les mêmes données avec le même $\lambda$, le même $K$ et
le même nombre d'itérations. L'article rapporte le même ordre sur son propre
instantané Amazon, et l'attribue à l'objectif sensible à la popularité plutôt qu'à
l'optimiseur - ce que confirme le point suivant.

\textbf{(ii) La pondération sensible à la popularité vaut à elle seule
$\PopGain$\,\% de HR@100.} La seule différence entre les deux lignes eALS du
Tableau~\ref{tab:quality} est $\alpha=0,4$ contre $\alpha=0$ ; tout le reste, y
compris $c_0$, la graine et le nombre d'itérations, est identique. La variante
uniforme se situe entre MLlib ALS et eALS complet, exactement là où l'argument de
l'article dit qu'elle devrait se trouver : une partie du gain vient d'une bonne
optimisation de l'objectif sur l'ensemble des données, une partie d'une bonne
pondération des données manquantes. Les deux modèles sont très au-dessus du plancher
MostPopular ($\PopFloor\times$ son HR@100).

\textbf{(iii) Sur la métrique qui compte en pratique - la précision pour un budget
de temps donné - eALS domine à chaque budget.} La Figure~\ref{fig:qualtime} est la
comparaison honnête, car une comparaison par itération masquerait le fait qu'une
itération ALS et une itération eALS ne représentent pas la même quantité de travail.
eALS atteint la meilleure précision qu'ALS n'atteint jamais (HR@100 $=\AlsBestHR$)
après $\EalsTimeToAls$\,s d'entraînement, contre $\AlsTimeToBest$\,s pour ALS, puis
continue de s'améliorer alors qu'ALS a plafonné : MLlib ALS est plat à partir de
l'itération~8, à HR@100 $\approx 0,118$. Notez que ceci n'est \emph{pas} une
affirmation sur la vitesse brute - à $K=32$ une itération ALS sur ces données est
moins chère que la nôtre, car MLlib est du code JVM natif et nous payons la
sérialisation Python. C'est une affirmation sur l'endroit où converge chaque
algorithme.
